\documentclass{article}
\usepackage{amsmath,amssymb,hyperref}
\newcommand{\Q}{\mathbb Q}\newcommand{\C}{\mathbb C}\newcommand{\Z}{\mathbb Z}
\begin{document}
\title{Targeted resources for the exceptional component-torsion gap}
\date{16 November 2026}
\maketitle
Background handoff only. The full exceptional Mathieu--Zhao cases remain open. No external reviewer was contacted, no new complete theorem was registered, and no finite phase/parameter campaign is proposed. The conditional comparison with safety relies on the existing internal local criterion, not independent verification of it.

\section{The latest actual control determines the search direction}
The agent's new \texttt{COHOMOLOGICAL\_PULLBACK\_AND\_A\_UNIFORM\_ORIGINAL\_MARKER.md} constructs actual invariant phases on both exceptional quotients with nonzero local components. Consequently exceptional invariance or perfectness alone cannot force torsion. Any closing proof must use the UNWEIGHTED eventual-purity identities. The pullback's fixed marker gives uniform interval moments, but does not identify the unweighted functional with those moments. The remaining universal target is
\[
 f\in\C[\mathrm{SL}_2/\Gamma],\quad \mathfrak h(f^m)=0\ (m\gg0)
 \quad\Longrightarrow\quad \nu_0(\sigma_f)=0,
 \qquad \Gamma=2O,2I.
\]

\section{A corrected component theorem yields an explicit valuation reduction}
The strongest directly matched new source is Halle--Nicaise, \emph{Neron models and base change}, \href{https://arxiv.org/abs/1209.5556}{arXiv:1209.5556}. References in this note use the cached 2012 preprint's numbering, not the later book's chapter numbering.

Chapter 4, Proposition 1.1.3 proves, for $0\to T\to G\to A\to0$ over a complete discretely valued field with algebraically closed residue field, that
\[
 \mathcal G^0(R)\longrightarrow\mathcal A^0(R)
 \text{ is onto},\qquad
 \Phi_T\longrightarrow\Phi_G\longrightarrow\Phi_A\longrightarrow0
 \text{ is exact}.
\]
Chapter 4, Proposition 3.3.1 additionally proves that the first map has torsion kernel. This is especially useful because Chapter 4, Section 1.1 corrects the perfect-residue-field case of Bosch--Xarles Lemma 4.2; the corrected result suffices for their subsequent component-group conclusions. One should use the corrected propositions rather than import that old lemma verbatim.

\textbf{Written deduction from these published propositions.} Put $K=\C((z))$, $R=\C[[z]]$. After a finite splitting cover let $T=\mathbf G_m^r$ be the ACTUAL toric subgroup of $G$, and let $s\in G(K)$ be the actual section point. Put $a=\pi(s)\in A(K)$. Choose a positive integer $n$ killing the finite component group $\Phi_A$. Then $[n]a\in\mathcal A^0(R)$. The published surjectivity gives a lift $b\in\mathcal G^0(R)$ of $[n]a$. In multiplicative notation set
\[
 t=[n]s\,b^{-1}\in T(K),\qquad
 v(s)=\frac1n\bigl(\operatorname{ord}_z\chi_i(t)\bigr)_{i=1}^r\in\Q^r,
\]
where the $\chi_i$ form a basis of $X^*(T)$. Then
\[
 \nu_G(s)=0\ \Longleftrightarrow\ v(s)=0.
\]
Moreover $v(s)$ is independent of $b$ and $n$ (for a fixed character basis and valuation normalization).

Proof: $\Phi_A$ is finite. The torsion-kernel proposition and $\Phi_T=\Z^r$ show that $\Phi_T\to\Phi_G$ is injective with finite cokernel. Hence it becomes an isomorphism over $\Q$. The component of $[n]s$ is the image of $\operatorname{val}_T(t)$ because $b$ has identity component. If $b'$ is another lift, $b'/b$ belongs to $T(K)\cap\mathcal G^0(R)$ and its toric valuation is zero by that injection. For two choices of $n$, take a common multiple and use the same independence of lifts. This also proves the criterion. The previously retained finite-cover reflection allows testing vanishing after a splitting cover. No claim of coordinate-independent numerical scaling between arbitrary covers is made.

This applies even when the abelian quotient has bad reduction and the semiabelian extension is nonsplit. It supplies an actual torus-valued correction, rather than fictitious independent torus coordinates on $G$. The characters are evaluated on $t\in T$, not on $s\in G$: they need not extend to characters of $G$. A scalar norm cannot replace the entire valuation vector.

For the canonical section, the retained local theorem would identify this valuation vector's vanishing with all-marker safety under purity. The NEW remaining constructive input is a usable description of $T$, the abelian component exponent, and a lift $b$ in the identity component. The source guarantees existence; this note does not implement an algorithm producing them for all phases. Nor does it prove purity makes these valuations zero.

\section{Logarithmic one-motives expose the gluing valuations, within scope}
Heer Zhao, \emph{Extending tamely ramified strict 1-motives into ket log 1-motives}, Forum Math. Sigma 9 (2021), e20, \href{https://doi.org/10.1017/fms.2021.5}{DOI 10.1017/fms.2021.5}, \href{https://arxiv.org/abs/1909.01046}{arXiv:1909.01046}, has the right type of explicit monodromy data.

In the published PDF, Theorem 3.1 extends a tamely ramified STRICT 1-motive. Section 4.1 defines its logarithmic monodromy pairing
\[
 \langle\ ,\ \rangle:Y\times X\longrightarrow
 \mathbf G_{m,\log}/\mathbf G_m.
\]
Proposition 4.1 identifies zero pairing with descent of the actual logarithmic section to an ordinary section. Proposition 4.2 subtracts the valuation part after choosing a uniformizer; Example 4.1 describes the good-reduction case. This is an avenue for expressing the distinguished section's obstruction from line-bundle trivializations and their orders, while retaining the Picard/NS/gluing maps.

The strictness hypothesis is essential: the introduction defines it as POTENTIALLY GOOD reduction of $G$. Residue characteristic zero removes wild ramification, but does not make a degenerating abelian variety potentially good. General genus degeneration is expressly allowed in the Mathieu theorem. Thus these propositions are not a ready-made universal vanishing theorem or an automatic general replacement for its Neron construction.

Raynaud, \emph{1-motifs et monodromie geometrique}, Asterisque 223 (1994), 295--319, \href{https://archive.numdam.org/item/AST_1994__223__295_0/}{original source}, is the underlying reference. For genuinely degenerating abelian quotients, Halle--Nicaise Chapter 4, Section 3.1 uses a rigid analytic uniformization; it explicitly warns that the resulting covering group is usually NOT algebraic. Applying the algebraic strict-one-motive theorem to that group without another comparison would be an error. The component-lift valuation deduction above avoids that issue entirely.

\section{A previously adopted moment theorem has a new conditional landing}
Francoise--Pakovich--Yomdin--Zhao, \emph{Moment vanishing problem and positivity: Some examples}, Bull. Sci. Math. 135 (2011), 10--32, \href{https://doi.org/10.1016/j.bulsci.2010.06.002}{DOI}, Proposition 3.3 and Theorem 3.3, prove more than a positivity statement: for a single polynomially parametrized interval and a polynomial complex density, nonzero total mass rules out an eventually vanishing power sequence of any nonzero polynomial on the parameterized curve. Their Theorem 4.1 also supplies the coprime-degree obstruction. Pakovich's \href{https://arxiv.org/abs/0910.2105}{0910.2105} supplies finite monodromy-relation descriptions for rational one-variable moment integrals (Theorem 3.2), and a strong irreducibility criterion in Theorem 4.2; it does not classify arbitrary higher-dimensional Haar periods.

These sources are ALREADY adopted in \texttt{REAL\_SEED\_DENSITY\_BOUNDS\_AND\_PERIOD\_TRACE.md}; they must not be advertised as newly discovered. A possible new landing is the current pullback family. If its UNWEIGHTED functional could be expressed exactly as
\[
 \mathfrak h(R(g_b))=\int_0^1 R(p(u))q(u)\,du
 \quad\text{for every polynomial }R,
\]
with $p$ nonconstant polynomial and $q$ polynomial, normalization gives $\int q=1$. The published theorem would then exclude EVERY nonzero eventually-pure polynomial $H(g_b)$ at once. Such a representation need not exist; the existing WEIGHTED uniform-marker identity cannot establish it. Atomic terms, varying curves, nonpolynomial densities or untracked branches invalidate this immediate application. An entire-family exclusion would settle this construction route, not the whole quotient.

\section{Protecting the search from plausible false shortcuts}
Markowsky--Phung, \href{https://arxiv.org/abs/2011.06344v2}{2011.06344v2}, Propositions 2--3, refute two stronger asymptotic guesses: power-moment growth need not equal the largest critical/endpoint value, and need not equal the maximum modulus on the integration interval. This does not refute the Muger--Tuset non-eventual-vanishing theorem used by the agent. It does rule out importing either of those larger maxima as an automatic singularity detector.

Results on zero loci of admissible normal functions, e.g. Brosnan--Pearlstein \href{https://arxiv.org/abs/math/0604345}{math/0604345} and \href{https://arxiv.org/abs/0803.3365}{0803.3365}, concern a different conclusion and do not establish torsion of this mixed toric component. The weight-minus-two torus is exactly what ordinary abelian arguments can miss. Likewise the logarithmic extension of a section is not itself a proof that its ordinary component vanishes.

A further abstract-level search lead was Estrada's 2020 nonperiodic real-analytic moment theorem, DOI \href{https://doi.org/10.37256/cm.132020307}{10.37256/cm.132020307}. The primary PDF fetch was blocked (403). It was not adopted or included as an actionable agent direction; full endpoint and periodicity hypotheses require reading before use.

\section{Recommended next obligation}
Keep the agent's full-target work. At a natural checkpoint, try to identify the toric correction $t=[n]\sigma\,b^{-1}$ in the actual Picard construction and ask whether UNWEIGHTED purity forces every valuation to vanish. This makes the remaining inference explicit without replacing it by symmetry, a scalar norm, an arbitrary local model, or another equivalent criterion presented as closure. No further marker search is needed on the existing pullback family.
\end{document}
